\documentclass[conference]{IEEEtran}

\usepackage{cite}
\usepackage{amsmath}
\usepackage{amssymb}
\usepackage{graphicx}
\usepackage{textcomp}
\usepackage{xcolor}
\usepackage{booktabs}
\usepackage{array}
\usepackage{tikz}
\usetikzlibrary{positioning,arrows.meta,shapes.geometric,fit}
\usepackage{url}
\usepackage[hidelinks]{hyperref}

\newcommand{\code}[1]{\path{#1}}

\begin{document}

\title{Detecting Fabricated Chat and Bank-Statement Screenshots:\\
Architecture, Implementation, and Evaluation of a Multi-Cue Forensic Pipeline}

\author{
\IEEEauthorblockN{Screenshot Forensics Project}
\IEEEauthorblockA{Technical Report \\ July 2026}
}

\maketitle

\begin{abstract}
Screenshots of chat conversations and bank statements are widely accepted as
informal proof of payment, identity, or communication, and are correspondingly
widely forged. This paper presents the design, implementation, and evaluation
of a working system that detects fabricated chat and bank-statement
screenshots by combining seven forensic cues -- Error Level Analysis (ELA),
JPEG Ghost double-compression analysis, JPEG quantization-table inspection,
block-wise noise-consistency analysis, offset-vector-voting copy-move
detection, EXIF/metadata inspection, and a novel optical-character-recognition
(OCR) based document-text cue -- into a single suspicion score. Cue scores
are combined by a RandomForest classifier trained on a real, publicly
available labeled document-forgery dataset, with a high-confidence override
layer that prevents a single conclusive signal (an editing-tool fingerprint,
a caught fake-document-generator watermark) from being diluted by noisier
pixel-level cues. We describe the full system architecture (a Flask backend
and React frontend, containerized with Docker), the end-to-end usage flow, and
the implementation of every cue in detail. We evaluate the system against both
a self-generated synthetic corpus and a real labeled receipt-forgery dataset,
report precision/recall/F1 results for a hand-weighted heuristic baseline
versus the trained classifier, and present a documented real-world case study
in which a fake bank statement was corrected from a false negative by the
override mechanism. We situate this work against 23 papers spanning classical
JPEG forensics, copy-move and splicing detection, deep-learning manipulation
detection, and document/identity-document forgery detection, and identify a
concrete literature gap: no prior published work specifically targets
digitally rendered chat or bank-statement screenshots as a content class
distinct from photographed or scanned documents. We report our results
honestly, including a documented recall ceiling and a training/deployment
domain gap, and outline concrete future work.
\end{abstract}

\begin{IEEEkeywords}
image forensics, document forgery detection, error level analysis, copy-move detection, JPEG compression forensics, optical character recognition, RandomForest, fraud detection
\end{IEEEkeywords}

\section{Introduction}
\label{sec:intro}

Digital screenshots occupy an unusual evidentiary position: they are trivially
easy to fabricate, yet routinely accepted at face value in disputes,
marketplace transactions, and informal loan or tenancy verification. Two
document classes account for the overwhelming majority of this misuse:
chat-application screenshots (used to fake proof of payment, e.g.\ peer-to-peer
transfer confirmations) and bank-statement screenshots or exports (used to
fake proof of funds or income). Industry reporting on ``fake payment
screenshot'' scams describes purpose-built generator tools that reproduce a
target payment app's fonts, logos, and layout closely enough that manual
visual inspection fails, and documents cases where victims released goods or
services against a screenshot alone before any funds actually
settled.\footnote{See, e.g., ``Fake Payment Screenshot Scams: How to Identify
\& Prevent UPI Fraud,'' Cashfree blog, and ``What Is the Fake Payment
Screenshot Scam? How Fraudsters Use Fake Receipts to Trick Retailers,''
LatestLY, both accessed 2026.}

This is, in forensic terms, a passive (blind) image authentication problem:
the verifier has only the image itself, with no reference original, no
digital signature, and -- critically for this content class -- no camera
sensor to fingerprint, since the ``photo'' is a rendered user-interface (UI)
screenshot rather than an optical capture. Decades of work in digital image
forensics address adjacent versions of this problem for photographs
(splicing, copy-move, resampling) and, more recently, for scanned or
photographed physical documents (receipts, ID cards, passports).
Section~\ref{sec:related} reviews this literature in depth. As
Section~\ref{sec:gap} makes explicit, however, none of it directly targets
screenshots of software UI as a content class, and this gap has concrete
consequences: several classical cues that are well validated on photographic
or scanned content are measurably weaker on flat, vector-rendered UI -- a
finding this paper documents empirically (Section~\ref{sec:results}) rather
than merely asserting.

\subsection{Contributions}
This paper makes the following contributions:
\begin{itemize}
\item A seven-cue forensic pipeline purpose-adapted for chat/bank-statement
screenshots, including a novel OCR-based document-text cue that targets what
turns out to be the most diagnostic signal for this content class: not pixel
statistics, but the document's own text.
\item A high-confidence override architecture that provably corrects a
documented real-world false negative (Section~\ref{sec:case-study}) without
increasing false-positive risk on the noisier pixel-level cues.
\item A trained classifier evaluated against a real, publicly available
labeled forgery dataset, with measured precision, recall, and F1 against both
a hand-weighted heuristic baseline and the trained model.
\item A literature review of 23 papers identifying where classical and
learned forensic techniques do and do not transfer to this content domain.
\end{itemize}

\subsection{Paper Organization}
Section~\ref{sec:related} reviews related work. Section~\ref{sec:architecture}
describes the system architecture. Section~\ref{sec:cues} details every
forensic cue's design and implementation.
Section~\ref{sec:classification} describes the decision logic.
Section~\ref{sec:implementation} covers implementation and deployment.
Section~\ref{sec:usage} walks through the end-to-end usage flow.
Section~\ref{sec:datasets} describes the datasets used.
Section~\ref{sec:results} reports experimental results.
Section~\ref{sec:discussion} discusses limitations, and
Section~\ref{sec:conclusion} concludes with future work.

\section{Related Work}
\label{sec:related}

This section reviews prior work across the subfields this project draws on.
Within each subsection we note not only what each paper contributes but,
where relevant, how well its assumptions transfer to the screenshot-forensics
setting, since several of the negative results in Section~\ref{sec:results}
trace directly to assumptions (photographic sensor noise, natural-image
statistics, an optical capture pipeline) that do not hold for a digitally
rendered UI.

\subsection{Classical JPEG and Compression-History Forensics}
The foundational insight exploited throughout this subfield is that JPEG
compression is lossy and content-dependent, so a region that has been edited
and re-saved carries a different compression history than the untouched
remainder of the image. Farid~\cite{farid2009ghosts} formalized this as
``JPEG Ghosts'': recompressing an image at a range of candidate qualities and
locating a region whose recompression error is minimized at a different
quality than the rest of the frame reveals a splice compressed at a different
generation. Lin~et~al.~\cite{lin2009fast} instead analyze the block-wise
histogram of discrete cosine transform (DCT) coefficients directly, since a
doubly JPEG-compressed region produces a distinctive periodic artifact that a
single compression does not. Bianchi and Piva contributed two closely related
detectors: a block-grained likelihood map for aligned and non-aligned double
JPEG compression~\cite{bianchi2012block}, and a companion method based on
integer periodicity maps for the non-aligned case~\cite{bianchi2012nonaligned},
which arises when a pasted region's block grid does not line up with the host
image's original $8\times8$ grid.

All four methods share an assumption this paper's results complicate: they
were validated on natural photographs or scanned documents, where legitimate
content is spatially smooth enough that compression artifacts concentrate
visibly at edited edges. A bank-statement screenshot is instead dense with
sharp, high-contrast text and hairline table borders throughout -- content
that itself produces strong, spatially widespread JPEG blocking artifacts
unrelated to any edit. Section~\ref{sec:results} reports that this content
mismatch, not an implementation defect, is the primary reason our JPEG-Ghost
and Error-Level-Analysis cues under-perform relative to these techniques'
reported accuracy on photographic benchmarks.

\subsection{Copy-Move Forgery Detection}
Copy-move -- duplicating a region within the same image, e.g.\ cloning a
background patch over an unwanted object, or in our threat model, duplicating
a transaction row to insert a fabricated one -- is one of the best-studied
forgery types because the duplicated content is pixel-identical (or
near-identical) to itself. Fridrich~et~al.~\cite{fridrich2003copymove}
proposed the original block-matching formulation: divide the image into
overlapping blocks, represent each with quantized DCT coefficients, and flag
blocks whose representations match at a spatial offset large enough to rule
out coincidental neighboring-block similarity. Popescu and
Farid~\cite{popescu2004exposing} refined the block representation using
principal component analysis (PCA) to improve robustness to minor noise or
compression differences. Huang~et~al.~\cite{huang2008sift} moved to
keypoint-based matching with SIFT features, trading some sensitivity to small
or smooth regions for strong robustness to rotation and rescaling.
Christlein~et~al.~\cite{christlein2012evaluation} later benchmarked fifteen
prominent feature sets under a common evaluation pipeline and found
keypoint-based methods (SIFT, SURF) and several block transforms (DCT, DWT,
PCA, Zernike moments) most robust overall, while cautioning that no single
method dominates across all post-processing scenarios -- a finding consistent
with our own decision (Section~\ref{sec:cue-copymove}) to combine block
hashing with an explicit geometric-consistency check rather than rely on raw
block matching alone.

Wu~et~al.'s BusterNet~\cite{wu2018busternet} represents the deep-learning-era
approach: a two-branch, end-to-end trainable network with one branch
localizing manipulation via visual artifacts and a second, dedicated branch
localizing copy-move regions via learned visual similarity, jointly
distinguishing source from target region. We did not adopt a learned
copy-move detector for this project; Section~\ref{sec:discussion} explains
why (principally, the absence of any labeled screenshot copy-move data to
train or validate one against) and instead describes a classical
offset-vector-voting detector in the Fridrich/Popescu lineage, extended with
a compactness gate motivated by a failure mode discovered empirically:
periodic decorative UI elements (repeated icons, evenly spaced list dividers)
produce coincidental block matches that a raw match count cannot distinguish
from a genuine rigid copy-paste.

\subsection{Splicing Detection and Sensor/Noise-Based Forensics}
Ng and Chang's Columbia dataset~\cite{ng2004columbia} -- 183 authentic images
and a matched set of spliced composites with no further post-processing --
was among the first public benchmarks built specifically for splicing
detection. Lyu, Pan, and Zhang~\cite{lyu2014exposing} address splicing
detection by exploiting the fact that image regions from different physical
sources typically carry different noise levels; their method estimates local
noise via the kurtosis of band-pass-filtered patches, then clusters regions
by estimated noise level to isolate a spliced region without needing a
reference image. Lukas, Fridrich, and Goljan~\cite{lukas2006digital} instead
exploit Photo Response Non-Uniformity (PRNU), a fixed noise pattern imprinted
by every camera's specific sensor, to attribute an image to its source camera
or flag a region whose PRNU does not match the rest of the frame.

This subfield illustrates a structural limitation for the present project
rather than a soft calibration issue: PRNU forensics is inapplicable to
screenshots outright, since a UI screenshot has no optical sensor and
therefore no PRNU signal to compare -- it is a rendering, not a photograph.
Our block-wise noise-consistency cue targets the same underlying idea as
Lyu~et~al.~\cite{lyu2014exposing} but must operate on whatever low-level
rendering/compression noise a screenshot happens to carry, which is a far
weaker signal than either PRNU or natural-image sensor noise.

\subsection{Deep-Learning-Era Manipulation Detection and Benchmark Datasets}
The CASIA Image Tampering Detection Evaluation Database~\cite{dong2013casia},
introduced by Dong, Wang, and Tan, provided a larger and more realistic
labeled tampering benchmark than the datasets available at the time and
remains widely used for training and evaluating both classical and learned
detectors. Zhou~et~al.~\cite{zhou2018learning} proposed a two-stream Faster
R-CNN (``RGB-N'') that fuses an RGB stream, which learns visual tampering
artifacts directly, with a dedicated noise stream built on steganalysis-style
rich filters, learning the same noise-inconsistency principle as
\cite{lyu2014exposing} end-to-end. Wang~et~al.~\cite{wang2020cnngenerated}
studied a related but distinct problem -- detecting wholly generated (not
just locally edited) CNN-synthesized images -- and showed a classifier
trained on a single generator family generalizes surprisingly well to unseen
architectures. Verdoliva's survey~\cite{verdoliva2020media} situates all of
the above within the field's broader shift from hand-crafted statistical
cues toward data-driven, end-to-end learned detectors, while noting that
data-driven methods inherit a generalization risk: performance can degrade
sharply on manipulation types or content domains absent from training data.

That generalization risk is precisely the argument this project's design
leans on in the opposite direction from most work in this subsection: rather
than train an end-to-end deep model with no labeled screenshot data to train
it on (Section~\ref{sec:gap}), we retain interpretable, individually
inspectable cues and train only a shallow classifier over their outputs, on
the closest available labeled data (Section~\ref{sec:datasets}).

\subsection{Document and Identity-Document Forgery Detection}
\label{sec:related-docs}
This subsection is the closest prior work to the present system's actual
deployment target. The SROIE dataset~\cite{huang2019icdar} (Huang~et~al.),
built for the ICDAR 2019 scanned-receipt OCR and information-extraction
competition, established a standard 1{,}000-receipt corpus with layout and
key-field annotations. Mart\'{i}nez Torn\'{e}s~et~al.'s ``Find it
again!''~dataset~\cite{martineztornes2023receipt} extends it directly for
forgery detection: 988 receipts drawn from SROIE, of which 163 carry
realistic forgeries (copy-paste, text imitation, deletion, and pixel-level
edits) with ground-truth region annotations. It is, to our knowledge, the
most directly applicable public labeled forgery dataset for this project's
problem, and Section~\ref{sec:datasets} describes its use here in detail.

Beyond receipts, a growing literature addresses identity-document forgery. A
preprint on deep-learning-based forgery attacks on document
images~\cite{anon2021forgery} examines both how deep networks can generate
convincing document forgeries and how the same technology can detect them. A
2026 systematic survey~\cite{anon2026survey} catalogs the shift in this area
from classical, hand-crafted approaches toward foundation-model-based
detectors, alongside purpose-built synthetic datasets created specifically
because real identity documents cannot be publicly released for privacy and
legal reasons. An ongoing competition series~\cite{anon2026competition} is
further evidence of how active this specific subfield now is.

The privacy constraint identified in~\cite{anon2026survey} for identity
documents applies with at least equal force to bank statements and
payment-app screenshots, which carry account numbers, balances, and personal
financial history -- and is precisely why this project
(Section~\ref{sec:datasets}) resorted to a self-generated synthetic corpus
for its own primary content domain, using the real receipt
dataset~\cite{martineztornes2023receipt} only as the closest available
substitute with genuine, non-synthetic ground truth.

\subsection{Adjacent Financial-Fraud-Detection Literature}
Two further pieces of literature, while not image forensics, directly
informed methodological choices in this project.
Hafemann~et~al.~\cite{hafemann2017learning} apply deep convolutional feature
learning to offline handwritten signature verification -- an older and
physically distinct form of financial-document authentication but one that
shares this project's core structure: an authenticator must distinguish
genuine from forged artifacts using only the static end artifact.
Jesus~et~al.'s Bank Account Fraud dataset suite~\cite{jesus2022turning}
established a widely used, realistic benchmark for fraud detection under
severe class imbalance; the practice it exemplifies -- training with
class-balanced weighting rather than naive accuracy optimization on a rare
positive class -- is the same principle behind this project's use of
\code{class_weight="balanced"} when training its RandomForest classifier
(Section~\ref{sec:classification}), since forged examples are a small
minority ($\approx$16\%) of the labeled data available.

\subsection{The Gap: Screenshot-Specific Forgery Detection}
\label{sec:gap}
Despite the breadth of work reviewed above -- spanning natural photographs
and, increasingly, scanned or photographed physical documents and identity
cards -- we found no academic literature or public dataset that specifically
targets chat-application or bank-statement screenshots: images that are a
direct digital rendering of software UI, never optically captured, and
typically saved once as a single-generation PNG or JPEG rather than scanned
or photographed. This is a meaningful gap for two reasons documented in
Sections~\ref{sec:results}--\ref{sec:discussion} of this paper. First,
several assumptions underlying the classical cues above do not hold for this
content, and our results show the corresponding cues transfer measurably
worse than their reported performance on photographic or scanned-document
benchmarks. Second, the practical harm is well documented outside the
academic literature (Section~\ref{sec:intro}). This project treats that gap
not as a reason to import ill-fitting photographic-forensics assumptions
unmodified, but as the reason to add a cue (Section~\ref{sec:cue-ocr}) that
targets what is actually diagnostic for this specific content class.

\section{System Architecture}
\label{sec:architecture}

The system is deployed as two containerized services orchestrated by Docker
Compose: a React single-page frontend served by nginx, and a Python/Flask
backend exposing a single analysis endpoint. Figure~\ref{fig:architecture}
shows the high-level component structure.

\begin{figure}[t]
\centering
\begin{tikzpicture}[
  box/.style={rectangle, draw, rounded corners, minimum width=6.6cm, minimum height=0.9cm, align=center, font=\small},
  cuebox/.style={rectangle, draw, minimum width=6.6cm, align=left, font=\scriptsize, inner sep=6pt},
  arr/.style={-{Stealth[length=2mm]}, thick}
]
\node[box] (client) {React Frontend (Vite, served via nginx)};
\node[box, below=0.55cm of client] (api) {Flask REST API: \code{POST /analyze}};
\node[cuebox, below=0.55cm of api] (cues) {
\textbf{Forensic Cue Extraction} (Section~\ref{sec:cues}):\\
1. Error Level Analysis \quad 2. JPEG Ghost\\
3. Quantization tables \quad 4. Noise consistency\\
5. Copy-move (offset voting) \quad 6. EXIF/metadata\\
7. Document text (OCR)
};
\node[box, below=0.55cm of cues] (classify) {Classification: RandomForest / heuristic (Sec.~\ref{sec:classification})};
\node[box, below=0.55cm of classify] (override) {High-Confidence Override Check};
\node[box, below=0.55cm of override] (response) {JSON Response: \code{fake_score}, \code{is_fake}, per-cue detail};

\draw[arr] (client) -- (api);
\draw[arr] (api) -- (cues);
\draw[arr] (cues) -- (classify);
\draw[arr] (classify) -- (override);
\draw[arr] (override) -- (response);
\draw[arr] (response.east) -- ++(1.3cm,0) |- (client.east);
\end{tikzpicture}
\caption{System architecture: request lifecycle for a single \code{/analyze} call.}
\label{fig:architecture}
\end{figure}

Images are processed entirely in memory: the backend never writes an
uploaded image to disk, logging only a SHA-256 hash for traceability. The
backend loads a trained RandomForest model (Section~\ref{sec:classification})
from \code{model_data/rf_model.joblib} at process start if present, and falls
back to a transparent weighted-average heuristic otherwise, so the service
remains functional in a minimal environment without scikit-learn or a
bundled model file.

\section{Forensic Cue Design}
\label{sec:cues}

Each cue returns a bounded suspicion score in $[0,1]$ together with the raw
evidence used to compute it, so the final verdict stays explainable rather
than a black-box output. Table~\ref{tab:cues} summarizes all seven cues and
the related work each draws on.

\begin{table}[t]
\centering
\caption{Forensic cues, their targets, and related work.}
\label{tab:cues}
\small
\begin{tabular}{@{}p{1.7cm}p{3.6cm}p{2.2cm}@{}}
\toprule
\textbf{Cue} & \textbf{Targets} & \textbf{Related work} \\
\midrule
ELA & Localized recompression artifacts & \cite{farid2009ghosts,bianchi2012block,bianchi2012nonaligned} \\
JPEG Ghost & Double compression at a different generation & \cite{farid2009ghosts,lin2009fast} \\
Quantization tables & Multiple compression generations & \cite{bianchi2012block,bianchi2012nonaligned} \\
Noise consistency & Locally inconsistent noise energy & \cite{lukas2006digital,lyu2014exposing} \\
Copy-move & Rigidly duplicated regions & \cite{fridrich2003copymove,popescu2004exposing,huang2008sift,christlein2012evaluation} \\
Metadata/EXIF & Editing-tool fingerprints, camera tags & (domain-specific) \\
Document text (OCR) & Generator watermarks, placeholder data & this work \\
\bottomrule
\end{tabular}
\end{table}

\subsection{Error Level Analysis}
Following Farid's recompression-diff principle~\cite{farid2009ghosts}, the
image is recompressed at a fixed JPEG quality and diffed against the input;
per-pixel error is aggregated into $16\times16$ blocks. Rather than flag
blocks whose error exceeds a single whole-image threshold -- which, on a
screenshot dense with text and table borders, flags most of the image
regardless of tampering -- each block is compared against a robust statistic
(median absolute deviation) of only its immediate spatial neighborhood. A
block that merely sits in a naturally detailed area is not flagged, since its
neighbors are similarly elevated; only a block that disagrees with its own
surroundings is. The size of the largest contiguous flagged cluster, not the
raw flagged-block count, drives the score, since a genuine localized edit
produces a compact cluster rather than scattered noise.

\subsection{JPEG Ghost}
This cue implements Farid's JPEG Ghost method~\cite{farid2009ghosts}
directly: the image is recompressed across ten candidate qualities
(50--95 in steps of 5) and, per block, the method looks for a genuine local
minimum in recompression error as a function of quality -- the signature of a
block's own native compression quality -- rather than naively taking the
single lowest-error quality across the whole range, which we found
empirically trivially selects the highest quality tested for almost any
content, producing no signal. Blocks are compared to their local
neighborhood rather than a single whole-image dominant quality, for the same
reason as Section~\ref{sec:cues}.A. Calibration testing
(Section~\ref{sec:results}) nonetheless found this cue unreliable on
screenshot content specifically -- in some cases scoring an untouched image
higher than its tampered counterpart -- so it is computed and shown for
transparency but deliberately excluded from the final weighted score.

\subsection{JPEG Quantization Tables}
A standard single JPEG encode carries at most two quantization tables (luma
and chroma); more than two on a re-opened file is a coarse signal of multiple
compression generations, in the spirit of Bianchi and Piva's double
compression detectors~\cite{bianchi2012block,bianchi2012nonaligned}, though
far less precise than their block-grained likelihood-map approach. In
practice this cue rarely fires, so it is retained as a cheap, low-weight
legacy signal.

\subsection{Block-Wise Noise Consistency}
A high-pass-filtered version of the image is divided into $32\times32$
blocks, and each block's local noise energy (standard deviation) is compared
to its spatial neighborhood using the same robust, local-outlier logic as
above, in the spirit of noise-inconsistency splicing
detection~\cite{lyu2014exposing} and PRNU-based source
consistency~\cite{lukas2006digital} -- adapted, since neither camera sensor
noise nor natural-image noise statistics are available for a rendered
screenshot, to whatever rendering/compression noise the image happens to
carry.

\subsection{Copy-Move Detection via Offset-Vector Voting}
\label{sec:cue-copymove}
Overlapping image blocks are perceptually hashed and grouped by hash; for
every pair of near-duplicate blocks that are spatially far apart, the
displacement vector between them is recorded as a ``vote,'' following the
block-matching lineage of~\cite{fridrich2003copymove,popescu2004exposing}. A
genuine rigid copy-paste produces many block pairs sharing the same
displacement vector, clustered in one spatially compact source area -- both
properties are required before the cue fires. This directly targets a
false-positive mode discovered empirically that a raw ``any two blocks
match'' count cannot distinguish: regularly spaced decorative UI elements
produce many block pairs that coincidentally share a common small offset
purely from the layout's own periodicity, without a spatially compact source
region, and are filtered out by the compactness requirement.

\subsection{EXIF / Metadata Analysis}
A genuine screenshot normally carries no EXIF data at all; this cue instead
flags its presence when specific tags are unusual for that fact: an
editing-tool \code{Software} tag (Photoshop, GIMP, etc.), camera
\code{Make}/\code{Model} tags, or GPS data, weighting an editing-tool
fingerprint especially heavily, since no legitimate screenshot-capture
pipeline stamps one. This is one of the two cues validated as reliable in
Section~\ref{sec:results}.

\subsection{Document Text Analysis (OCR)}
\label{sec:cue-ocr}
Motivated directly by the gap identified in Section~\ref{sec:gap} and the
case study in Section~\ref{sec:case-study}, this cue runs Tesseract OCR over
the image and checks the extracted text against two pattern families: known
fake-document-generator brand markers and disclaimer phrasing (e.g.\ a
generator site's own watermark, or ``specimen'' / ``not a real document''
boilerplate), and placeholder data a generator defaults to when unedited
(generic names such as ``John Doe,'' and obviously sequential or
repeated-digit account/routing numbers). Unlike every pixel-level cue above,
this cue reads what the document \emph{says} rather than how its bytes were
compressed, and is the other cue validated as reliable in
Section~\ref{sec:results}.

\section{Classification and the High-Confidence Override}
\label{sec:classification}

Cue scores are combined into a single \code{fake_score} by a
\code{RandomForestClassifier} (300 trees, max depth 6,
\code{class_weight="balanced"} per the class-imbalance precedent
in~\cite{jesus2022turning}) trained on labeled data
(Section~\ref{sec:datasets}). A pure weighted average -- the system's
original design and still its fallback when no trained model is available --
has a documented failure mode this project encountered directly: a genuinely
conclusive single cue (an editing-tool EXIF tag, an OCR-caught generator
watermark) can be diluted below the decision threshold by several other,
noisier cues that happen to read elevated on legitimate content, as
Section~\ref{sec:case-study}'s case study shows concretely.

To close this gap without simply re-weighting the noisy cues higher (which
would raise false positives on legitimate dense documents instead), an
explicit override rule sits on top of either classifier:
\begin{equation}
\label{eq:override}
\text{is\_fake} =
\begin{cases}
\text{true} & \text{if } s_{\text{composite}} \geq \tau \\
\text{true} & \text{if } \exists\, c \in \mathcal{O} : s_c \geq \theta_c \\
\text{false} & \text{otherwise}
\end{cases}
\end{equation}
where $s_{\text{composite}}$ is the classifier's output, $\tau$ is the
decision threshold (0.5), $\mathcal{O} = \{\text{metadata\_score},
\text{text\_score}\}$ is the set of override-eligible cues, and
$\theta_c = 0.8$ for each. Only cues independently validated as
high-precision, low-false-positive signals belong in $\mathcal{O}$; the
composite score $s_{\text{composite}}$ is still reported unmodified alongside
an explicit \code{override_reason} field, so an override is never applied
silently.

\section{Implementation}
\label{sec:implementation}

\subsection{Backend}
The backend is implemented in Python using Flask, with Pillow for image
decoding/encoding, NumPy for the numerical cue computations described in
Section~\ref{sec:cues}, pytesseract/Tesseract for the OCR cue
(Section~\ref{sec:cue-ocr}), and scikit-learn/joblib for the trained
classifier (Section~\ref{sec:classification}). A single route,
\code{POST /analyze}, accepts a multipart image upload (PNG, JPEG, or WebP,
capped at 10\,MB), validates the file extension and non-emptiness, and
returns a JSON response containing the composite verdict, the raw per-cue
feature vector, per-cue diagnostic detail (e.g.\ dominant JPEG quality,
matched OCR flags, clone-detection offset votes), and a base64-encoded ELA
heatmap image for display. A companion \code{GET /status} route serves as a
health check. Errors are handled explicitly: an unreadable file returns
HTTP~400 with a descriptive message, and any unexpected exception during
analysis is logged server-side and returns a generic HTTP~500 without leaking
internals to the client.

\subsection{Frontend}
The frontend is a React single-page application built with Vite. It handles
drag-and-drop or click-to-browse file selection with client-side validation
(file type, size), displays a live preview before analysis, and renders the
full result set after a successful \code{/analyze} call: a verdict card with
an icon and a horizontal severity meter for the composite score, a
side-by-side original/ELA-heatmap comparison, and a per-cue breakdown list in
which every cue is rendered as its own severity meter with an explanatory
note drawn from that cue's diagnostic detail. A visible banner explains when
the high-confidence override (Section~\ref{sec:classification}) rather than
the composite score determined the verdict. A ``Download report'' action
exports the full JSON response as a file for external record-keeping.

\subsection{Deployment and Testing}
Both services are containerized with dedicated Dockerfiles and orchestrated
via Docker Compose, with nginx reverse-proxying frontend API calls to the
Flask container. The backend test suite (pytest) covers each cue function in
isolation, end-to-end API behavior, the override mechanism, and an
explicit accuracy-regression test tied to the case study in
Section~\ref{sec:case-study}. The frontend test suite (Vitest,
React Testing Library) covers upload validation behavior. A calibration
tool, \code{tools/calibrate.py}, generates the synthetic corpus described in
Section~\ref{sec:datasets} and reports per-cue win-rates against it; a
companion \code{tools/train_model.py} retrains the bundled classifier from
any dataset following the same labeled-image-directory convention as the
real dataset in Section~\ref{sec:datasets}.

\section{Usage Flow}
\label{sec:usage}

Figure~\ref{fig:usage} shows the end-to-end user interaction sequence, which
maps directly onto the implementation described in
Section~\ref{sec:implementation}.

\begin{figure}[t]
\centering
\begin{tikzpicture}[
  step/.style={rectangle, draw, rounded corners, minimum width=6.4cm, minimum height=0.8cm, align=center, font=\small},
  arr/.style={-{Stealth[length=2mm]}, thick}
]
\node[step] (s1) {1. User drags in or selects an image file};
\node[step, below=0.4cm of s1] (s2) {2. Client-side validation (type, size $\leq$10\,MB)};
\node[step, below=0.4cm of s2] (s3) {3. Preview shown; user clicks ``Analyze''};
\node[step, below=0.4cm of s3] (s4) {4. \code{POST /analyze} (multipart upload)};
\node[step, below=0.4cm of s4] (s5) {5. Backend runs all 7 cues (Sec.~\ref{sec:cues})};
\node[step, below=0.4cm of s5] (s6) {6. Classifier + override determine verdict};
\node[step, below=0.4cm of s6] (s7) {7. JSON response rendered: verdict, meters, cue detail};
\node[step, below=0.4cm of s7] (s8) {8. User downloads report or starts a new scan};

\draw[arr] (s1) -- (s2);
\draw[arr] (s2) -- (s3);
\draw[arr] (s3) -- (s4);
\draw[arr] (s4) -- (s5);
\draw[arr] (s5) -- (s6);
\draw[arr] (s6) -- (s7);
\draw[arr] (s7) -- (s8);
\end{tikzpicture}
\caption{End-to-end usage flow, from image selection to report export.}
\label{fig:usage}
\end{figure}

If validation fails at step 2 (unsupported type or oversized file), the user
sees an inline error and no request is sent. If the backend cannot decode the
uploaded bytes as an image, step 4 returns HTTP~400 and the client surfaces
that message directly rather than a generic failure. A loading state (spinner
and disabled controls) covers the interval between steps 4 and 7, which
completes in well under one second for a typical screenshot-sized image in
our testing.

\section{Datasets}
\label{sec:datasets}

\subsection{Synthetic Corpus}
In the absence of any public labeled screenshot-forgery dataset
(Section~\ref{sec:gap}), a synthetic corpus generator
(\code{tools/calibrate.py}) produces paired authentic/tampered chat- and
bank-statement-style images across four tamper types: (i) a textured region
edited and recompressed at a different JPEG quality before being pasted back
and the whole image re-saved -- the classical scenario
Section~\ref{sec:cues}.A--B targets; (ii) copy-move duplication of a textured
block; (iii) a low-texture ``sloppy'' flat overwrite, the hardest case for
compression-artifact cues; and (iv) an editing-tool EXIF fingerprint left on
an otherwise untouched image. This corpus's paired structure (same
underlying content, with vs. without the tamper) makes it possible to
measure whether a cue's score reliably moves in the correct direction for a
specific tamper type, rather than merely whether real and fake images differ
on average -- a distinction Section~\ref{sec:results} shows can be misleading
when the two classes also differ in unrelated ways (e.g.\ an extra
compression generation).

\subsection{Real Dataset: ``Find it again!''~Receipt Forgery Dataset}
The trained classifier (Section~\ref{sec:classification}) uses Mart\'{i}nez
Torn\'{e}s~et~al.'s dataset~\cite{martineztornes2023receipt}: 988 scanned
receipt images derived from SROIE~\cite{huang2019icdar}, split into 577
train / 192 validation / 218 test examples, of which 163 total ($\approx$16\%)
carry realistic forgeries with ground-truth region annotations. This is real
labeled forgery data, unlike the synthetic corpus above, but it is scanned
paper receipts rather than digitally rendered screenshots -- a domain gap
discussed candidly in Section~\ref{sec:discussion}. The training script
extracts the same seven cue scores the served API computes for every image
and trains the RandomForest directly on them.

\section{Experimental Results}
\label{sec:results}

\subsection{Synthetic-Corpus Calibration: Per-Cue Reliability}
Rather than compare score distributions across unrelated real/fake samples --
which conflates the tamper's effect with incidental differences in content
and compression history between samples -- we measured, for each of 25
paired authentic/tampered images, how often the tampered version scored
higher than its own authentic pair (``win rate''). Table~\ref{tab:calibration}
reports this per cue.

\begin{table}[t]
\centering
\caption{Paired win-rate per cue across four synthetic tamper types ($n=25$).}
\label{tab:calibration}
\small
\begin{tabular}{@{}lc@{}}
\toprule
\textbf{Cue} & \textbf{Win rate (tampered $>$ authentic)} \\
\midrule
ELA & 6 / 25 \\
JPEG Ghost & 11 / 25 \\
Quantization tables & 0 / 25 \\
Noise consistency & 8 / 25 \\
Copy-move & 7 / 25 \\
Metadata & 6 / 25 (100\% within EXIF-tamper subset) \\
\bottomrule
\end{tabular}
\end{table}

The low win-rates for ELA, JPEG Ghost, and noise consistency -- well below
what the literature in Section~\ref{sec:related} reports for photographic
content -- are the empirical basis for the content-domain-mismatch argument
made throughout this paper, and directly motivated (a) excluding JPEG Ghost
from the final score and (b) weighting metadata and document-text far more
heavily than the pixel-level cues in the trained classifier.

\subsection{Real-Dataset Results: Heuristic vs. Trained Classifier}
On the ``Find it again!'' test split (218 receipts, 35 forged), the original
hand-weighted heuristic combiner was compared against the trained
RandomForest classifier. Table~\ref{tab:results} reports the results.

\begin{table}[t]
\centering
\caption{Held-out test-set performance ($n=218$, 35 forged).}
\label{tab:results}
\small
\begin{tabular}{@{}lcc@{}}
\toprule
\textbf{Metric} & \textbf{Heuristic} & \textbf{RandomForest} \\
\midrule
Precision & 42.3\% & \textbf{100\%} (0 false pos.) \\
Recall & 31.4\% & 28.6\% \\
F1 & 0.361 & \textbf{0.444} \\
Accuracy & 82.1\% & 88.5\% \\
\bottomrule
\end{tabular}
\end{table}

The trained model's zero false positives is a meaningful result given the
framing in Section~\ref{sec:intro}: in a workflow where a flagged document
routes to manual review, a high-precision, lower-recall classifier minimizes
wasted investigator time and the reputational risk of wrongly challenging a
genuine customer's document, at the cost of still missing roughly 71\% of the
real forgeries in this test set. Neither figure should be read as a claim of
solved detection; Section~\ref{sec:discussion} discusses this recall ceiling
directly.

\subsection{Case Study: A Real-World False Negative and Its Fix}
\label{sec:case-study}
During manual testing, a bank-statement screenshot generated by a known
fake-statement-generator site -- carrying that site's own watermark and
placeholder customer data (``John Doe,'' account number ``123456789'') --
scored $\text{fake\_score}=0.235$ (``Likely Authentic'') from the composite
classifier alone, despite the document-text cue
(Section~\ref{sec:cue-ocr}) itself scoring $1.0$ on the same image. The
composite score under-weighted this single conclusive signal because two
moderately elevated but individually inconclusive pixel-level cues (ELA and
noise consistency, both plausibly explained by the document's own text
density rather than tampering) pulled the blended score down. Applying the
high-confidence override (Eq.~\ref{eq:override}) corrected the verdict to
``Possible Fake'' with $\text{override\_reason}=\text{text\_score}$, without
needing to re-weight the unreliable pixel cues. This single case is the
concrete motivation for the override architecture and is retained here as a
worked example rather than a synthetic illustration.

\subsection{Cross-Domain Validation}
Because the trained classifier was trained on scanned receipts
(Section~\ref{sec:datasets}) but deployed against chat/bank-statement
screenshots, its own predicted probability was checked directly against this
project's screenshot test images rather than assumed to transfer. On the
watermarked bank-statement screenshot from Section~\ref{sec:case-study}, the
trained model's own probability was 0.235 -- the override, not the model's
cross-domain transfer, is what corrected the verdict. This is reported
plainly as a limitation rather than elided: the model measurably does not
fully transfer across the domain gap identified in
Section~\ref{sec:related-docs}, which is precisely why the override layer is
a permanent part of the architecture rather than a stopgap.

\section{Discussion and Limitations}
\label{sec:discussion}

\begin{itemize}
\item \textbf{Domain mismatch, not implementation error}, explains most
pixel-cue weakness. Sections~\ref{sec:related} and \ref{sec:results} together
show that ELA, JPEG Ghost, and noise consistency all under-perform their
literature-reported accuracy specifically on flat, text-dense, vector-rendered
UI content -- exactly the content type absent from the photographic and
scanned-document benchmarks those techniques were developed against.
\item \textbf{The recall ceiling (28.6\% on real data) is a genuine,
unsolved limitation}, not merely a conservative design choice: roughly seven
in ten real forgeries in the held-out test set are missed by the trained
classifier's own signal, before the override layer (which only helps when
metadata or document-text independently fires).
\item \textbf{No deep-learning component.} Section~\ref{sec:related}
documents the field's shift toward end-to-end learned detectors precisely
because they outperform hand-crafted cues when sufficient labeled data
exists; this project deliberately did not pursue that path, given the
absence of labeled screenshot-forgery data large enough to train one without
high overfitting risk.
\item \textbf{Training/deployment domain gap.} The only real, non-synthetic
labeled forgery data available is scanned receipts, not screenshots -- a
direct consequence of the same privacy/legal constraint on real
financial-document data that the identity-document literature documents for
its own adjacent domain.
\item \textbf{Copy-move false-positive risk remains non-zero} for dense,
regularly repeated UI; the compactness gate narrows but does not eliminate
this risk, consistent with Christlein~et~al.'s finding~\cite{christlein2012evaluation}
that no single copy-move method dominates every scenario.
\item \textbf{OCR reliability depends on image resolution and legibility}; a
blurry photograph of a printed statement may not OCR cleanly enough to catch
a watermark that a clean digital screenshot would.
\end{itemize}

\section{Conclusion and Future Work}
\label{sec:conclusion}

This paper presented a forensic pipeline for a document class -- chat and
bank-statement screenshots -- that the academic literature does not directly
address, despite extensive and mature work on adjacent problems. The system
combines seven cues spanning classical JPEG/compression forensics, copy-move
detection, metadata analysis, and a novel OCR-based document-text cue,
combined via a RandomForest trained on the closest available real labeled
data and protected by a high-confidence override layer motivated by a
genuine documented detection failure. Measured results show real progress --
zero false positives and improved F1 over the original hand-weighted
heuristic on held-out real data -- alongside equally real, openly reported
limitations: weak pixel-level cue transfer to this content domain, a recall
ceiling under 30\%, and an unresolved training/deployment domain gap.

Three directions follow directly from the gaps this paper documents. First,
constructing a larger, more varied synthetic screenshot-forgery corpus that
more closely matches the receipt dataset's diversity of tamper types could
narrow the domain gap without requiring real, privacy-sensitive screenshot
data. Second, the deep-learning direction reviewed in
Section~\ref{sec:related} -- in particular a noise-stream architecture in the
spirit of RGB-N~\cite{zhou2018learning} -- becomes viable once such a corpus
exists at sufficient scale. Third, given that this project's single most
reliable cue reads document text rather than pixels, extending it beyond
watermark/placeholder pattern-matching toward semantic consistency checks
(e.g.\ does the stated bank's name and address consistently appear across the
document) is a promising, comparatively low-cost next step directly
motivated by the case study in Section~\ref{sec:case-study}.

\bibliographystyle{IEEEtran}
\bibliography{references}

\end{document}
